\documentclass[a4paper]{article}
\usepackage{amsmath}
\usepackage{amsfonts}
\usepackage{amssymb}
\usepackage{graphicx}
\begin{document}

\noindent \large Auteur: Siebe Haché \\
\noindent \large Studierichting: Informatica \\
\noindent \large Oefeningenreeks 3F nr. 2.5 \\

\medskip
\normalsize
Bepaal de schuine asymptoot door een euclidische deling te maken: \\

$ f(x) = \dfrac{\sqrt{2}x^2 + \sqrt{3}x - \sqrt{2}}{x + \sqrt{6}} $\\

\bigskip

\begin{tabular}{rccc|l}
     & $+\sqrt{2}x^2$  & $+\sqrt{3}x$  & $-\sqrt{2}$  & $x + \sqrt{6}$ \\ \cline{5-5} 
$\_$ & $+\sqrt{2}x^2$  & $+2\sqrt{3}x$  &              & $\sqrt{2}x - \sqrt{3}$ \\ \cline{2-3}
     &         & $-\sqrt{3}x$   & $-\sqrt{2}$        &\\
     & $\_$     & $-\sqrt{3}x$  & $-3\sqrt{2}$       &\\ \cline{3-4}
     &          &       &  $+2\sqrt{2}$             & \\ 
\end{tabular}\\

Schuine asymptoot:\\

$A: y = \sqrt{2}x - \sqrt{3}$ \\

$ f(x) - A = \dfrac{2\sqrt{2}}{x + \sqrt{6}}$\\

\begin{tabular}{c|ccc}
$x$		   &     & $-\sqrt{6}$ &     \\ \hline
$f(x) - A$ & $-$ & $|$         & $+$ \\
\end{tabular}\\

Dus:\\

$ \lim\limits_{x \rightarrow +\infty} (f(x) - A) = 0^+ \Rightarrow x \rightarrow +\infty : f > A $ \\ 

$ \lim\limits_{x \rightarrow -\infty} (f(x) - A) = 0^- \Rightarrow x \rightarrow -\infty : f < A $ \\ 

\begin{thebibliography}{999}
%\bibitem{a} Referentie
\end{thebibliography}
\end{document}